\textbf{Open-vocabulary object search.} Zero-shot semantic navigation scores where to explore with vision-language models: VLFM builds a language-grounded value map over frontiers~\cite{yokoyama2023vlfm}, OneMap keeps a reusable open-vocabulary feature map for multi-object search~\cite{busch2024one}, and HOV-SG organises open-vocabulary features into a floor, room and object scene graph for language-grounded navigation~\cite{werby2024hierarchical}. HM3D-OVON broadens object-goal navigation to goals given in free-form language~\cite{yokoyama2024hm}. The perception side relies on open-vocabulary detectors such as OWL-ViT~\cite{minderer2022simple}. Our detector is a simulated stand-in for this text-similarity scoring, and our SPL-style metric follows the standard evaluation recommendations for embodied navigation~\cite{anderson2018evaluation}.

\textbf{Multi-robot exploration and search.} Frontier-based exploration and its multi-robot variants remain the reference family, benchmarked for example in Explore-Bench~\cite{xu2022explore}. For semantic search, GoalSwarm coordinates several aerial robots on a shared top-down semantic map~\cite{james2026goalswarm}, and COMRES-VLM lets a vision-language model coordinate frontier assignment over shared maps~\cite{wang2025comres}. These systems share maps or state directly; the message format is not the object of study.

\textbf{Language between heterogeneous robots.} Cladera et al.\ present air-ground collaboration for language-specified missions, where a ground and an aerial robot share information under intermittent communication~\cite{cladera2025air}. MHRC uses language models for decentralised collaboration among heterogeneous robots~\cite{yu2024mhrc}; Chen et al.\ compare centralised and decentralised language-model dialogue for multi-robot planning with attention to token budgets~\cite{chen2023scalable}; CoELA builds embodied agents that plan and communicate in natural language when communication is costly~\cite{zhang2023building}; and Choe et al.\ let heterogeneous robots request and offer help in natural language~\cite{choe2025ask}. These works evaluate complete language stacks. They do not hold the agents fixed and swap only the message format.

\textbf{What to communicate.} In multi-agent reinforcement learning, a survey organises communication by what is sent, to whom and under which constraints~\cite{zhu2022survey}, and Hill et al.\ compare an engineered intention message with an end-to-end learned protocol~\cite{hill2025communicating}. Our study is in this spirit but with hand-written channels: we compare two engineered formats and separate message content from the way a location is referred to.

\textbf{Gap.} To our knowledge the closest systems~\cite{cladera2025air,yu2024mhrc} report end-to-end results for language-coordinated heterogeneous teams. We did not find a controlled comparison of natural-language against fixed-schema messages for ground-aerial open-vocabulary search at matched budget on paired layouts; this toy study addresses the information-format part of that question only.
